\section{Experimental Setup}
\subsection{Dataset}
We evaluate on LoCoMo~\cg{maharana2024evaluating}, a long-horizon conversational memory benchmark containing 10 multi-session dialogues with 1{,}540 non-adversarial QA pairs across four question categories: single-hop (282), multi-hop (321), temporal (96), and open-domain (841). We exclude the adversarial category following the Memory-R1 evaluation protocol~\cg{yan2025memory}. Following Memory-R1, we use a 1:1:8 train/validation/test split, yielding 152 training, 81 validation, and 1{,}307 test QA pairs. LoCoMo is well matched to our setting: it requires persistent memory over many turns, supports both single-hop and multi-hop reasoning, and includes temporal and open-domain questions that test different aspects of memory access.

\subsection{Model and Hardware}
All trainable components use Qwen2.5-3B-Instruct as the base policy model. Experiments are run on a single NVIDIA H200 GPU with 16 CPU cores. This setup supports single-node staged GRPO training and fixed-answer-agent evaluation.

We use GRPO as the default policy-optimization method, following Search-R1 and Memory-R1~\cg{jin2025search, yan2025memory}. For the Memory Manager and Retrieve Agent, prompts are formatted as one-shot JSON-generation problems with rule-based reward functions. The Answer Agent is trained to produce short answers conditioned on retrieved memories. During Memory Manager and Retrieve Agent training, the Answer Agent is frozen and used only to compute reward.

\subsection{Baselines}
We organize baselines into three categories.

\paragraph{Context-based.}
\textbf{Full Context} feeds the entire conversation history into the base Qwen2.5-3B-Instruct model without explicit memory or retrieval, relying on the context window alone.
\textbf{Oracle} uses the gold evidence annotations provided in LoCoMo, representing an upper bound on retrieval quality.

\paragraph{Retrieval-based.}
\textbf{RAG (BM25)} applies BM25 keyword retrieval over a flat memory bank of extracted facts.
\textbf{RAG (Embedding)} uses text-embedding-3 dense retrieval over the same flat memory bank.
\textbf{Semantic XPath}~\cg{semanticxpath2026} generates deterministic XPath-style queries over oracle tree-structured memory.

\paragraph{RL-based.}
\textbf{Memory R1}~\cg{yan2025memory} is a flat-memory system with an RL-trained Memory Manager and Answer Agent, trained on the same base model and the same data split for fair comparison.\footnote{Memory-R1 does not provide a public codebase. We built a publicly available implementation at \url{https://github.com/YifanLiu2/Struct_Memory_R1.git}.}

\subsection{Evaluation Metrics}
We report three answer-quality metrics. \textbf{F1} is the token-level F1 score between the predicted and gold answers, measuring lexical overlap. \textbf{B1} is BLEU-1, a unigram precision metric that captures answer conciseness alongside correctness. \textbf{J} is the LLM-as-judge accuracy, where GPT-4o mini evaluates whether the predicted answer is semantically correct given the gold answer. All three metrics are reported per question category (multi-hop, single-hop, temporal, open-domain) and as a weighted overall average reflecting the natural category distribution in LoCoMo.
